Existen varios aspectos de un alimento que podemos tomar en cuenta para clasificarlo. Estos son:

\begin{itemize}
    \item Función
        \begin{itemize}
	        \item Energéticos: Grasas y carbohidratos
	        \item De construcción: Vitaminas y minerales
	    \end{itemize}
    
    \item Origen
        \begin{itemize}
	        \item Animal
	        \item Vegetal
        \end{itemize}
    
    \item Durabilidad
        \begin{itemize}
	        \item Perecedero
	        \item No perecedero
        \end{itemize}
    
    \item Grado de procesamiento
        \begin{itemize}
	        \item Natural o minimamente procesado
	        \item Ingrediente culinario
	        \item Procesado
	        \item Ultraprocesado
        \end{itemize}
    
    \item Grupo alimenticio\footnote{Basado en el plato del buen comer}
        \begin{itemize}
	        \item Frutas y verduras
	        \item Cereales y tubérculos
	        \item Leguminosas y alimentos de origen animal
        \end{itemize}
        
\end{itemize}

\section{Tipos de fibras}



\begin{itemize}
    \item Dietarias
        \begin{itemize}
	        \item No son descompuestas por el metabolismo humano, por lo que unicamente pasan por el intestino sin ser procesadas.
	        \item Ayudan a la digestión.
        \end{itemize}
        
    \item Ácido-resistentes
        \begin{itemize}
	        \item Son descompuestas por el ácido gástrico.
	        \item Sos componentes son absorbidos y utilizados por el cuerpo.
        \end{itemize}
\end{itemize}

\section{Tipos de análisis alimentarios}



\begin{itemize}
    \item Bromatológico: En este se realizan pruebas de una serie amplia de parámetros. De estos estudios se obtiene la información para la tabla nutrimental del etiquetado de alimentos.
    \item Proximal: Es un análisis general de unicamente ciertos parámetros. Los parámetros medibles:
        \begin{itemize}
		    \item Agua
		    \item Proteína
		    \item Carbohidratos
		    \item Grasa
		    \item Ceniza
	    \end{itemize}
    \item Sanitario: Son análisis que buscan presencia de patógenos y/o deterioro, esto para prevenir enfermedades transmitidas por alimentos (ETAs').
\end{itemize}